\documentclass[11pt,a4paper]{article}
\input{../../sablony/pripravy/priprava.sty}
\usepackage{tikz,pgfplots}
\pgfplotsset{compat=1.18}
\nastavPripravu{1. ročník}{Kvadratické rovnice a nerovnice}{119}{Slovní úlohy na kvadratické rovnice I}
\begin{document}
\blok{Slovní úlohy na kvadratické rovnice I}
Geometrie a rozměry: obsah, šířka, obvod, fyzikální smysl kořenů.
\blok{Metoda a návaznost na graf}
\begin{itemize}
\item Zaveď neznámou, sestav vztah a po vyřešení zkontroluj kladnost délek a původní zadání.
\item Fyzikálně nemožný algebraický kořen musí být výslovně odmítnut.
\end{itemize}
\blok{Společně řešené příklady}
\begin{prikladReseniBox}[title={Příklad 1}]
Obdélník má obsah $60\,\mathrm{cm^2}$ a delší strana je o $7\,\mathrm{cm}$ větší. Určete délky stran.
\tcblower
\[
\begin{aligned}
x(x+7)&=60\\[2pt]
x^2+7x-60&=0\\[2pt]
(x-5)(x+12)&=0\\[2pt]
x&=5\,\mathrm{cm}\quad (x=-12\text{ nelze})\\[2pt]
\text{Strany: }&5\,\mathrm{cm};\ 12\,\mathrm{cm}
\end{aligned}
\]
\end{prikladReseniBox}
\begin{prikladReseniBox}[title={Příklad 2}]
Čtvercový pozemek má po rozšíření každé strany o $3\,\mathrm m$ obsah $100\,\mathrm{m^2}$. Jaká byla původní strana?
\tcblower
\[
\begin{aligned}
(x+3)^2&=100\\[2pt]
x+3&=10\quad (x>0)\\[2pt]
x&=7\,\mathrm m
\end{aligned}
\]
\end{prikladReseniBox}
\begin{prikladReseniBox}[title={Příklad 3}]
Model výšky míčku je $h(t)=-5t^2+20t+1$ (metry, sekundy). Kdy je míček ve výšce $16\,\mathrm m$?
\tcblower
\[
\begin{aligned}
-5t^2+20t+1&=16\\[2pt]
t^2-4t+3&=0\\[2pt]
(t-1)(t-3)&=0\\[2pt]
t&=1\,\mathrm s\ \text{nebo}\ 3\,\mathrm s
\end{aligned}
\]
\end{prikladReseniBox}
\blok{Samostatné procvičování}
\begin{zadaniReseniBox}[title={Samostatná práce 1}]
Dráha brzdění je v modelu $s=\frac{v^2}{2a}$, $a=5\,\mathrm{m/s^2}$. Jakou počáteční rychlost (v m/s) odpovídá dráha $40\,\mathrm m$?
\tcblower
\[
\begin{aligned}
40&=\frac{v^2}{10}\\[2pt]
v^2&=400\\[2pt]
v&=20\,\mathrm{m/s}\quad(v\ge0)
\end{aligned}
\]
\end{zadaniReseniBox}
\begin{zadaniReseniBox}[title={Samostatná práce 2}]
Výnos výrobního procesu je modelován $Z(q)=-2q^2+40q-150$ (tis. Kč). Pro které nezáporné $q$ je $Z(q)>0$?
\tcblower
\[
\begin{aligned}
-2q^2+40q-150&>0\\[2pt]
q^2-20q+75&<0\\[2pt]
(q-5)(q-15)&<0\\[2pt]
K&=(5;15)
\end{aligned}
\]
\end{zadaniReseniBox}

\blok{Geometrická aplikace: zlatý řez}
\begin{definiceBox}[title={Definice - Zlatý řez}]
\textbf{Zlatý řez} je rozdělení úsečky délky $a+b$ ($a>b>0$),
při kterém poměr celé úsečky k delší části odpovídá poměru delší části ke kratší:
\[\frac{a+b}{a}=\frac ab=\varphi.\]
Číslo $\varphi$ se nazývá \textbf{zlatý poměr}.
\end{definiceBox}
\begin{center}
\begin{tikzpicture}[x=0.95cm,y=0.95cm]
\draw[thick] (0,0)--(8.09,0);
\draw[thick] (0,-.16)--(0,.16) (5,-.16)--(5,.16) (8.09,-.16)--(8.09,.16);
\node[above] at (2.5,.12) {$a$}; \node[above] at (6.55,.12) {$b$};
\node[below] at (0,-.15) {$A$};\node[below] at (5,-.15) {$B$};\node[below] at (8.09,-.15) {$C$};
\end{tikzpicture}
\end{center}
Proč vznikne právě kvadratická rovnice? Z $a/b=\varphi$ plyne
$b=a/\varphi$. Dále $\frac{a+b}{a}=1+\frac b a=1+1/\varphi$.
Proto
\[
\begin{aligned}
\varphi&=1+\frac1\varphi\quad(\varphi>1),\\
\varphi^2-\varphi-1&=0,\\
D&=(-1)^2-4\cdot1\cdot(-1)=5,\\
\varphi&=\frac{1+\sqrt5}{2}\doteq1{,}618.
\end{aligned}
\]
Záporný kořen $(1-\sqrt5)/2$ nelze použít jako poměr kladných délek.
\textbf{Geometrie:} Zlatý poměr se objevuje například v poměru úhlopříčky a strany
\emph{pravidelného pětiúhelníku}. U \emph{zlatého obdélníku} po odříznutí čtverce
zůstane menší obdélník s týmž poměrem stran. Tento princip je známý už
z antické geometrie a práce s podobností trojúhelníků.
\poznamka{Při hodině rozlište matematickou přesnost od popularizačních tvrzení: nelze automaticky označovat libovolný obraz, stavbu či přírodní útvar za projev zlatého řezu. Dobrovolně: Je-li kratší část 8 cm, spočítejte delší část a délku celé úsečky.}

\blok{Volitelná rezerva}
\poznamka{Studenti si vymyslí obdobný příklad se změnou jednoho koeficientu a vysvětlí, jak se změní počet kořenů či řešení.}
\blok{Domácí úkol}
\begin{zadaniReseniBox}[title={Domácí úkol}]
\begin{enumerate}[label=\alph*),itemsep=2pt]
\item Čtvercový pozemek má po rozšíření každé strany o $3\,\mathrm m$ obsah $100\,\mathrm{m^2}$. Jaká byla původní strana?
\item Model výšky míčku je $h(t)=-5t^2+20t+1$ (metry, sekundy). Kdy je míček ve výšce $16\,\mathrm m$?
\item Dráha brzdění je v modelu $s=\frac{v^2}{2a}$, $a=5\,\mathrm{m/s^2}$. Jakou počáteční rychlost (v m/s) odpovídá dráha $40\,\mathrm m$?
\item Výnos výrobního procesu je modelován $Z(q)=-2q^2+40q-150$ (tis. Kč). Pro které nezáporné $q$ je $Z(q)>0$?
\item Obdélníková obrazovka má obsah $96\,\mathrm{cm^2}$, šířka je o $4\,\mathrm{cm}$ menší než výška. Určete rozměry.
\item Model výkonové ztráty odporu je $P=I^2R$, $R=8\,\Omega$. Pro jaký nezáporný proud dosahuje $P=72\,\mathrm W$?
\item Časový model letu dronu je $h(t)=-2t^2+12t$ metrů. Kdy je ve výšce $10\,\mathrm m$?
\item Obdélník má obsah $60\,\mathrm{cm^2}$ a delší strana je o $7\,\mathrm{cm}$ větší. Určete délky stran.
\end{enumerate}
\tcblower
\begin{enumerate}[label=\alph*),itemsep=2pt]
\item
\[\begin{aligned}
(x+3)^2&=100\\[2pt]
x+3&=10\quad (x>0)\\[2pt]
x&=7\,\mathrm m
\end{aligned}\]
\item
\[\begin{aligned}
-5t^2+20t+1&=16\\[2pt]
t^2-4t+3&=0\\[2pt]
(t-1)(t-3)&=0\\[2pt]
t&=1\,\mathrm s\ \text{nebo}\ 3\,\mathrm s
\end{aligned}\]
\item
\[\begin{aligned}
40&=\frac{v^2}{10}\\[2pt]
v^2&=400\\[2pt]
v&=20\,\mathrm{m/s}\quad(v\ge0)
\end{aligned}\]
\item
\[\begin{aligned}
-2q^2+40q-150&>0\\[2pt]
q^2-20q+75&<0\\[2pt]
(q-5)(q-15)&<0\\[2pt]
K&=(5;15)
\end{aligned}\]
\item
\[\begin{aligned}
x(x-4)&=96\\[2pt]
x^2-4x-96&=0\\[2pt]
(x-12)(x+8)&=0\\[2pt]
\text{Výška}&=12\,\mathrm{cm},\ \text{šířka}=8\,\mathrm{cm}
\end{aligned}\]
\item
\[\begin{aligned}
72&=8I^2\\[2pt]
I^2&=9\\[2pt]
I&=3\,\mathrm A\quad(I\ge0)
\end{aligned}\]
\item
\[\begin{aligned}
-2t^2+12t&=10\\[2pt]
t^2-6t+5&=0\\[2pt]
(t-1)(t-5)&=0\\[2pt]
t&\in\{1;5\}\,\mathrm s
\end{aligned}\]
\item
\[\begin{aligned}
x(x+7)&=60\\[2pt]
x^2+7x-60&=0\\[2pt]
(x-5)(x+12)&=0\\[2pt]
x&=5\,\mathrm{cm}\quad (x=-12\text{ nelze})\\[2pt]
\text{Strany: }&5\,\mathrm{cm};\ 12\,\mathrm{cm}
\end{aligned}\]
\end{enumerate}
\end{zadaniReseniBox}
\poznamka{Orientační čas: úvod 5 min, společné příklady 18 min, samostatně 15 min, kontrola a shrnutí 7 min. Rezerva je dobrovolná. Číselné parametry modelových aplikací nejsou skutečně naměřená data.}
\end{document}
